\section{Conclusion}
\label{sec:conclusion}

We have presented ResearchOS, a framework that transforms unstructured scientific papers into composable \emph{research skills}---structured primitives with defined functional interfaces---and organizes them into a Research Skill Graph with typed relations capturing prerequisite, enhancement, substitution, conflict, composition, and refinement dependencies. A constraint-aware composition engine leverages this graph to generate novel method designs through goal decomposition, candidate retrieval, directed acyclic graph construction, gap-bridging, and iterative self-critique, producing three levels of progressively actionable output: Method Design Cards, Research Proposals, and Executable Experiment Plans.

Our experiments yield four key findings. First, structured composition over a skill graph significantly outperforms direct generation, retrieval-augmented generation, and unstructured skill lists, demonstrating that explicit relational structure is essential for high-quality method design. Second, graph relations---particularly conflict and substitution edges---are critical for producing feasible designs; removing them leads to compositions with internal contradictions. Third, gap-bridging is the primary driver of genuine novelty, synthesizing connective mechanisms that go beyond recombination of existing techniques. Fourth, future-grounded evaluation via time-split validation confirms that ResearchOS designs anticipate research directions found in subsequently published papers, with Method Hit Rates substantially above both baselines and random chance.

Looking ahead, we see three directions for extending this work: broadening skill extraction to additional research domains beyond LLM reasoning, closing the loop between design generation and automated execution \cite{lu2024aiscientist}, and developing community-maintained skill repositories that grow and improve as the research literature evolves.
